\lecture{8}{16. Marts 2026}{Linear isotropic materials and the uniaxial stress}

\begin{problem}{4.2}
  For an isotropic material, obtain the relationship
  \begin{equation*}
    \epsilon_{ij} = \frac{1}{2\mu} \sigma_{ij} - \frac{\lambda}{2\mu \left( 2\mu + 3\lambda \right)} \delta_{ij} \sigma_{kk}
  \end{equation*}
  from the expression
  \begin{equation*}
    \sigma_{ij} = 2\mu \epsilon_{ij} + \lambda\delta_{ij}\epsilon_{kk}
  .\end{equation*}
\end{problem}

\begin{problem}{4.3a}
  Show that the principal axes of stress correspond with the principal axes of strain for an isotropic, linearly elastic material.
\end{problem}

\begin{problem}{4.5}
  For steel, $E = \num{30e6}$ and $G = \num{12e6}$ ($\text{force}/\text{length}^2$).
  The components of strain at a point within this material are given by
  \begin{equation*}
    \Vec{\epsilon} = \begin{bmatrix}
      \num{0,004}  & \num{0,001}  & 0\\
    \num{0,001}  & \num{0,006}  & \num{0,004} \\
    0 & \num{0,004}  & \num{0,001} \\
    \end{bmatrix}
  .\end{equation*}
  Compute the corresponding components of the stress tensor.
\end{problem}
